\chapter{Estado del arte}
\label{cap2:arte}

La gestión eficiente y el dimensionamiento preventivo de las infraestructuras de telecomunicaciones móviles dependen de la capacidad de anticipar la demanda de tráfico de red para garantizar los estándares de calidad de servicio.

Para alcanzar este objetivo se utiliza la predicción de tráfico de red modelado como serie temporal. Entre las primeras aproximaciones al problema destacan los métodos estadísticos clásicos basados en supuestos de estacionariedad como los modelos autorregresivos y de medias móviles \cite{arma,arma1,arma2,arma3,arima1,arima2,sarima1}. Son enfoques sencillos y ligeros computacionalmente pero con grandes ineficiencias debido a la elevada cantidad de suposiciones de linealidad y estacionariedad que no concuerdan con las características del tráfico de red actual.

La introducción y desarrollo de la inteligencia artificial y, en concreto, los modelos de aprendizaje profundo supuso una ampliación en las capacidades de predicción de las series temporales. El tráfico de red ha sido procesado por una gran variedad de modelos desde arquitecturas sencillas como los MLPs hasta estructuras más complejas como las neuronas recurrentes \cite{SKIMAcasaseca,tfgalicia}.

En un contexto de crecimiento exponencial de las fuentes de tráfico de red y gran facilidad de acceso a dispositivos de telecomunicaciones por parte del público general, el tráfico de red es cada vez más complejo con periodicidades multidimensionales, no linealidades, correlaciones espaciales y cambios en su evolución temporal. La necesidad de realizar una ingeniería de características con intervención humana para ajustar los modelos a las series temporales analizadas supone una inversión cada vez mayor \cite{deeplearning}. Por ello, aunque los modelos neuronales superan sistemáticamente la precisión obtenida por los métodos autorregresivos en la predicción de series temporales \cite{tfgalicia},los modelos de aprendizaje automático que no involucran \textit{deep learning} han quedado obsoletos.

Los modelos de aprendizaje profundo son la siguiente evolución en la predicción de series temporales introduciendo múltiples capas neuronales que permiten la extracción automática de las características del tráfico de red.

Las series temporales que modelan el tráfico de red pueden estar relacionadas entre sí porque cada serie pertenece a un punto geográfico concreto que puede ser espacialmente cercano a otro punto que genere una serie temporal distinta. Los modelos puramente recurrentes tratan los flujos de tráfico como series temporales independientes ignorando las correlaciones espaciales existentes en el tráfico de red \cite{lstm,lstm2}.

Parte del tráfico de red es generado por dispositivos móviles capaces de desplazarse entre diferentes posiciones geográficas, afectando al tráfico medido en cada una. Dichas posiciones pueden pertenecer a una misma área como el casco urbano de una ciudad. El desarrollo de las reces convolucionales LSTM permitió el procesado de las correlaciones espaciales entre series temporales provenientes de puntos geográficos cercanos \cite{shi2015}.

En esta línea se han realizado trabajos previos de predicción de series temporales correlacionadas a través de un plano geográfico urbano utilizando redes convolucionales y recurrentes \cite{shi2015,zhang2018,zhou2022,shen2023,he2024}.

Este TFG se centra en la predicción de tráfico de red utilizando los datos de tráfico de la ciudad de Milán \cite{milandataset}. Sobre este dataset se han realizado varios estudios proponiendo formas más eficientes de retropropagación en régimen incremental \cite{SKIMAcasaseca} y diferentes variantes de arquitecturas de redes neuronales convolucionales y recurrentes \cite{zhang2018,shen2023,he2024,tfgalicia}. 

La capacidad de movimiento de los dispositivos generadores de tráfico añade cambios en el patrón subyacente del tráfico de red que degradan la capacidad predictiva de los modelos \textit{offline}. Para superar esta vulnerabilidad se han desarrollado los modelos \textit{online}, capaces de adaptarse dinámicamente a los datos procesados fuera del periodo de entrenamiento \cite{conceptdrift,conceptdrift1,conceptdrift2}.

Los trabajos mencionados anteriormente validan el estudio del tráfico de red en entornos urbanos como una matriz bidimensional de series temporales a través de modelos neuronales profundos que combinan capas convolucionales y recurrentes. Sin embargo, no implementan ninguna capacidad de adaptación frente a cambios en el patrón subyacente del tráfico de red \cite{shi2015,zhang2018,zhou2022,shen2023,he2024,tfgalicia}.

Este TFG plantea la hipótesis de que la hibridación entre modelos neuronales convolucionales para extraer los patrones geográficos y temporales para realizar predicciones suponen una mejora respecto a las arquitecturas anteriormente empleadas utilizando una porción mayor del mismo dataset. También busca evaluar si estas mismas arquitecturas se benefician del entrenamiento en régimen \textit{online} frente al entrenamiento clásico \textit{offline}.